\IfFileExists{stacks-project.cls}{%
\documentclass{stacks-project}
}{%
\documentclass{amsart}
}

% For dealing with references we use the comment environment
\usepackage{verbatim}
\newenvironment{reference}{\comment}{\endcomment}
%\newenvironment{reference}{}{}
\newenvironment{slogan}{\comment}{\endcomment}
\newenvironment{history}{\comment}{\endcomment}

% For commutative diagrams we use Xy-pic
\usepackage[all]{xy}

% We use 2cell for 2-commutative diagrams.
\xyoption{2cell}
\UseAllTwocells

% We use multicol for the list of chapters between chapters
\usepackage{multicol}

% This is generally recommended for better output
\usepackage{lmodern}
\usepackage[T1]{fontenc}

% For cross-file-references
\usepackage{xr-hyper}

% Package for hypertext links:
\usepackage{hyperref}

% For any local file, say "hello.tex" you want to link to please
% use \externaldocument[hello-]{hello}
\externaldocument[introduction-]{introduction}
\externaldocument[conventions-]{conventions}
\externaldocument[sets-]{sets}
\externaldocument[categories-]{categories}
\externaldocument[topology-]{topology}
\externaldocument[sheaves-]{sheaves}
\externaldocument[sites-]{sites}
\externaldocument[stacks-]{stacks}
\externaldocument[fields-]{fields}
\externaldocument[algebra-]{algebra}
\externaldocument[brauer-]{brauer}
\externaldocument[homology-]{homology}
\externaldocument[derived-]{derived}
\externaldocument[simplicial-]{simplicial}
\externaldocument[more-algebra-]{more-algebra}
\externaldocument[smoothing-]{smoothing}
\externaldocument[modules-]{modules}
\externaldocument[sites-modules-]{sites-modules}
\externaldocument[injectives-]{injectives}
\externaldocument[cohomology-]{cohomology}
\externaldocument[sites-cohomology-]{sites-cohomology}
\externaldocument[dga-]{dga}
\externaldocument[dpa-]{dpa}
\externaldocument[sdga-]{sdga}
\externaldocument[hypercovering-]{hypercovering}
\externaldocument[schemes-]{schemes}
\externaldocument[constructions-]{constructions}
\externaldocument[properties-]{properties}
\externaldocument[morphisms-]{morphisms}
\externaldocument[coherent-]{coherent}
\externaldocument[divisors-]{divisors}
\externaldocument[limits-]{limits}
\externaldocument[varieties-]{varieties}
\externaldocument[topologies-]{topologies}
\externaldocument[descent-]{descent}
\externaldocument[perfect-]{perfect}
\externaldocument[more-morphisms-]{more-morphisms}
\externaldocument[flat-]{flat}
\externaldocument[groupoids-]{groupoids}
\externaldocument[more-groupoids-]{more-groupoids}
\externaldocument[etale-]{etale}
\externaldocument[chow-]{chow}
\externaldocument[intersection-]{intersection}
\externaldocument[pic-]{pic}
\externaldocument[weil-]{weil}
\externaldocument[adequate-]{adequate}
\externaldocument[dualizing-]{dualizing}
\externaldocument[duality-]{duality}
\externaldocument[discriminant-]{discriminant}
\externaldocument[derham-]{derham}
\externaldocument[local-cohomology-]{local-cohomology}
\externaldocument[algebraization-]{algebraization}
\externaldocument[curves-]{curves}
\externaldocument[resolve-]{resolve}
\externaldocument[models-]{models}
\externaldocument[functors-]{functors}
\externaldocument[equiv-]{equiv}
\externaldocument[pione-]{pione}
\externaldocument[etale-cohomology-]{etale-cohomology}
\externaldocument[proetale-]{proetale}
\externaldocument[relative-cycles-]{relative-cycles}
\externaldocument[more-etale-]{more-etale}
\externaldocument[trace-]{trace}
\externaldocument[crystalline-]{crystalline}
\externaldocument[spaces-]{spaces}
\externaldocument[spaces-properties-]{spaces-properties}
\externaldocument[spaces-morphisms-]{spaces-morphisms}
\externaldocument[decent-spaces-]{decent-spaces}
\externaldocument[spaces-cohomology-]{spaces-cohomology}
\externaldocument[spaces-limits-]{spaces-limits}
\externaldocument[spaces-divisors-]{spaces-divisors}
\externaldocument[spaces-over-fields-]{spaces-over-fields}
\externaldocument[spaces-topologies-]{spaces-topologies}
\externaldocument[spaces-descent-]{spaces-descent}
\externaldocument[spaces-perfect-]{spaces-perfect}
\externaldocument[spaces-more-morphisms-]{spaces-more-morphisms}
\externaldocument[spaces-flat-]{spaces-flat}
\externaldocument[spaces-groupoids-]{spaces-groupoids}
\externaldocument[spaces-more-groupoids-]{spaces-more-groupoids}
\externaldocument[bootstrap-]{bootstrap}
\externaldocument[spaces-pushouts-]{spaces-pushouts}
\externaldocument[spaces-chow-]{spaces-chow}
\externaldocument[groupoids-quotients-]{groupoids-quotients}
\externaldocument[spaces-more-cohomology-]{spaces-more-cohomology}
\externaldocument[spaces-simplicial-]{spaces-simplicial}
\externaldocument[spaces-duality-]{spaces-duality}
\externaldocument[formal-spaces-]{formal-spaces}
\externaldocument[restricted-]{restricted}
\externaldocument[spaces-resolve-]{spaces-resolve}
\externaldocument[formal-defos-]{formal-defos}
\externaldocument[defos-]{defos}
\externaldocument[cotangent-]{cotangent}
\externaldocument[examples-defos-]{examples-defos}
\externaldocument[algebraic-]{algebraic}
\externaldocument[examples-stacks-]{examples-stacks}
\externaldocument[stacks-sheaves-]{stacks-sheaves}
\externaldocument[criteria-]{criteria}
\externaldocument[artin-]{artin}
\externaldocument[quot-]{quot}
\externaldocument[stacks-properties-]{stacks-properties}
\externaldocument[stacks-morphisms-]{stacks-morphisms}
\externaldocument[stacks-limits-]{stacks-limits}
\externaldocument[stacks-cohomology-]{stacks-cohomology}
\externaldocument[stacks-perfect-]{stacks-perfect}
\externaldocument[stacks-introduction-]{stacks-introduction}
\externaldocument[stacks-more-morphisms-]{stacks-more-morphisms}
\externaldocument[stacks-geometry-]{stacks-geometry}
\externaldocument[moduli-]{moduli}
\externaldocument[moduli-curves-]{moduli-curves}
\externaldocument[examples-]{examples}
\externaldocument[exercises-]{exercises}
\externaldocument[guide-]{guide}
\externaldocument[desirables-]{desirables}
\externaldocument[coding-]{coding}
\externaldocument[obsolete-]{obsolete}
\externaldocument[fdl-]{fdl}
\externaldocument[index-]{index}

% Theorem environments.
%
\theoremstyle{plain}
\newtheorem{theorem}[subsection]{Theorem}
\newtheorem{proposition}[subsection]{Proposition}
\newtheorem{lemma}[subsection]{Lemma}

\theoremstyle{definition}
\newtheorem{definition}[subsection]{Definition}
\newtheorem{example}[subsection]{Example}
\newtheorem{exercise}[subsection]{Exercise}
\newtheorem{situation}[subsection]{Situation}

\theoremstyle{remark}
\newtheorem{remark}[subsection]{Remark}
\newtheorem{remarks}[subsection]{Remarks}

\numberwithin{equation}{subsection}

% Macros
%
\def\lim{\mathop{\mathrm{lim}}\nolimits}
\def\colim{\mathop{\mathrm{colim}}\nolimits}
\def\Spec{\mathop{\mathrm{Spec}}}
\def\Hom{\mathop{\mathrm{Hom}}\nolimits}
\def\Ext{\mathop{\mathrm{Ext}}\nolimits}
\def\SheafHom{\mathop{\mathcal{H}\!\mathit{om}}\nolimits}
\def\SheafExt{\mathop{\mathcal{E}\!\mathit{xt}}\nolimits}
\def\Sch{\mathit{Sch}}
\def\Mor{\mathop{\mathrm{Mor}}\nolimits}
\def\Ob{\mathop{\mathrm{Ob}}\nolimits}
\def\Sh{\mathop{\mathit{Sh}}\nolimits}
\def\NL{\mathop{N\!L}\nolimits}
\def\CH{\mathop{\mathrm{CH}}\nolimits}
\def\proetale{{pro\text{-}\acute{e}tale}}
\def\etale{{\acute{e}tale}}
\def\QCoh{\mathit{QCoh}}
\def\Ker{\mathop{\mathrm{Ker}}}
\def\Im{\mathop{\mathrm{Im}}}
\def\Coker{\mathop{\mathrm{Coker}}}
\def\Coim{\mathop{\mathrm{Coim}}}

% Boxtimes
%
\DeclareMathSymbol{\boxtimes}{\mathbin}{AMSa}{"02}

%
% Macros for moduli stacks/spaces
%
\def\QCohstack{\mathcal{QC}\!\mathit{oh}}
\def\Cohstack{\mathcal{C}\!\mathit{oh}}
\def\Spacesstack{\mathcal{S}\!\mathit{paces}}
\def\Quotfunctor{\mathrm{Quot}}
\def\Hilbfunctor{\mathrm{Hilb}}
\def\Curvesstack{\mathcal{C}\!\mathit{urves}}
\def\Polarizedstack{\mathcal{P}\!\mathit{olarized}}
\def\Complexesstack{\mathcal{C}\!\mathit{omplexes}}
% \Pic is the operator that assigns to X its picard group, usage \Pic(X)
% \Picardstack_{X/B} denotes the Picard stack of X over B
% \Picardfunctor_{X/B} denotes the Picard functor of X over B
\def\Pic{\mathop{\mathrm{Pic}}\nolimits}
\def\Picardstack{\mathcal{P}\!\mathit{ic}}
\def\Picardfunctor{\mathrm{Pic}}
\def\Deformationcategory{\mathcal{D}\!\mathit{ef}}
\usepackage{cleveref}
\crefname{originallemma}{original lemma}{original lemmas}
\Crefname{originallemma}{Original lemma}{Original lemmas}
\makeatletter
\newlabel{etale-cohomology-lemma-category-loc-cst-is-colimit}{{0GL2}{source}{}{sourcepointer-1}{}}
\newlabel{etale-cohomology-lemma-category-loc-cst-is-colimit@cref}{{[originallemma][0][]0GL2}{[1][0][]source}{}{}{}}
\makeatother
\crefname{originallemma}{original lemma}{original lemmas}
\Crefname{originallemma}{Original lemma}{Original lemmas}
\makeatletter
\newlabel{etale-cohomology-lemma-characterize-locally-constant-module}{{0GKC}{source}{}{sourcepointer-2}{}}
\newlabel{etale-cohomology-lemma-characterize-locally-constant-module@cref}{{[originallemma][0][]0GKC}{[1][0][]source}{}{}{}}
\makeatother
\crefname{originallemma}{original lemma}{original lemmas}
\Crefname{originallemma}{Original lemma}{Original lemmas}
\makeatletter
\newlabel{etale-cohomology-lemma-sp-isom-proper-torsion-loc-ac}{{0GJW}{source}{}{sourcepointer-3}{}}
\newlabel{etale-cohomology-lemma-sp-isom-proper-torsion-loc-ac@cref}{{[originallemma][0][]0GJW}{[1][0][]source}{}{}{}}
\makeatother
\crefname{originallemma}{original lemma}{original lemmas}
\Crefname{originallemma}{Original lemma}{Original lemmas}
\makeatletter
\newlabel{limits-lemma-descend-dimension-d}{{0EY2}{source}{}{sourcepointer-4}{}}
\newlabel{limits-lemma-descend-dimension-d@cref}{{[originallemma][0][]0EY2}{[1][0][]source}{}{}{}}
\makeatother
\crefname{originallemma}{original lemma}{original lemmas}
\Crefname{originallemma}{Original lemma}{Original lemmas}
\makeatletter
\newlabel{limits-lemma-descend-finite-presentation}{{01ZM}{source}{}{sourcepointer-5}{}}
\newlabel{limits-lemma-descend-finite-presentation@cref}{{[originallemma][0][]01ZM}{[1][0][]source}{}{}{}}
\makeatother
\crefname{originallemma}{original lemma}{original lemmas}
\Crefname{originallemma}{Original lemma}{Original lemmas}
\makeatletter
\newlabel{limits-lemma-descend-smooth}{{0C0C}{source}{}{sourcepointer-6}{}}
\newlabel{limits-lemma-descend-smooth@cref}{{[originallemma][0][]0C0C}{[1][0][]source}{}{}{}}
\makeatother
\crefname{originallemma}{original lemma}{original lemmas}
\Crefname{originallemma}{Original lemma}{Original lemmas}
\makeatletter
\newlabel{limits-lemma-eventually-proper}{{081F}{source}{}{sourcepointer-7}{}}
\newlabel{limits-lemma-eventually-proper@cref}{{[originallemma][0][]081F}{[1][0][]source}{}{}{}}
\makeatother
\crefname{originaltheorem}{original theorem}{original theorems}
\Crefname{originaltheorem}{Original theorem}{Original theorems}
\makeatletter
\newlabel{theorem-constructible-shriek}{{0GL0}{source}{}{sourcepointer-8}{}}
\newlabel{theorem-constructible-shriek@cref}{{[originaltheorem][0][]0GL0}{[1][0][]source}{}{}{}}
\makeatother
\title{Proper smooth pushforwards of uniformly torsion locally constant sheaves are finite locally constant}
\author{The Stacks Project Authors}
\date{stacks-project:more-etale.tex:local-snapshot-20261012; GNU FDL 1.2 or later}
\begin{document}
\maketitle
\setcounter{subsection}{1}
\noindent\textit{Editorial note.} All delivered mathematical text is the author's own, copied from the pinned Stacks Project chapter file more-etale.tex; no statement or proof step was rewritten, reordered or removed. The counted span is the author environment and its own proof environment. References inside the span are handed to the assembler as navigation stubs that carry the canonical tag of the cited tag only; whether each of them is load-bearing is recorded on the stub itself.\par\medskip
\noindent\textit{Allowed standard prerequisites (editorial).}\begin{itemize}
\item Constructibility of proper pushforwards (Theorem theorem-constructible-shriek, canonical tag 0GL0). cited at proof lines 5064-5065. In the final Noetherian n-torsion case it provides constructibility of the sheaves R\textasciicircum{}if\_*F, the input to the specialization and characterization lemmas.
\item Locally constant sheaves as a colimit (Etale Cohomology, Lemma etale-cohomology-lemma-category-loc-cst-is-colimit, canonical tag 0GL2). cited at proof lines 5051-5055. Lets the locally constant sheaf F descend to a finite type locally constant F\_i on some X\_i in the Noetherian descent.
\item Locally constant characterization (Etale Cohomology, Lemma etale-cohomology-lemma-characterize-locally-constant-module, canonical tag 0GKC). cited at proof lines 5069-5070. The terminal criterion: constructible sheaves whose specialization maps are isomorphisms are locally constant.
\end{itemize}\medskip
\ifdefined\hypertarget\hypertarget{sourcepointer-1}{}\fi
\noindent\textit{Source reference.} Original lemma 0GL2: sources/stacks-chow/tags.txt, source lines 20707--20707. Cited at proof lines 5051-5055: F descends to a finite type locally constant sheaf F\_i on some X\_i.\par\medskip
\ifdefined\hypertarget\hypertarget{sourcepointer-2}{}\fi
\noindent\textit{Source reference.} Original lemma 0GKC: sources/stacks-chow/tags.txt, source lines 20682--20682. Cited at proof lines 5069-5070: constructible sheaves whose specialization maps are isomorphisms are locally constant, finishing the proof.\par\medskip
\ifdefined\hypertarget\hypertarget{sourcepointer-3}{}\fi
\noindent\textit{Source reference.} Original lemma 0GJW: sources/stacks-chow/tags.txt, source lines 20666--20666. Cited at proof lines 5066-5068: the specialization maps of R\textasciicircum{}if\_*F are isomorphisms.\par\medskip
\ifdefined\hypertarget\hypertarget{sourcepointer-4}{}\fi
\noindent\textit{Source reference.} Original lemma 0EY2: sources/stacks-chow/tags.txt, source lines 18677--18677. Cited at proof lines 5046-5050: one of the three Limits lemmas used to improve f\_i after increasing i.\par\medskip
\ifdefined\hypertarget\hypertarget{sourcepointer-5}{}\fi
\noindent\textit{Source reference.} Original lemma 01ZM: sources/stacks-chow/tags.txt, source lines 2599--2599. Cited at proof lines 5042-5045: descend f to a finitely presented morphism f\_i : X\_i -> S\_i over a finite type Z[1/n]-subalgebra.\par\medskip
\ifdefined\hypertarget\hypertarget{sourcepointer-6}{}\fi
\noindent\textit{Source reference.} Original lemma 0C0C: sources/stacks-chow/tags.txt, source lines 15082--15082. Cited at proof lines 5046-5050: after increasing i, f\_i is smooth.\par\medskip
\ifdefined\hypertarget\hypertarget{sourcepointer-7}{}\fi
\noindent\textit{Source reference.} Original lemma 081F: sources/stacks-chow/tags.txt, source lines 10215--10215. Cited at proof lines 5046-5050: after increasing i, f\_i is proper.\par\medskip
\ifdefined\hypertarget\hypertarget{sourcepointer-8}{}\fi
\noindent\textit{Source reference.} Original theorem 0GL0: sources/stacks-chow/tags.txt, source lines 20705--20705. Cited at proof lines 5064-5065: the sheaves R\textasciicircum{}if\_*F are constructible in the Noetherian n-torsion case.\par\medskip
\begin{proposition}
\label{proposition-loc-cst-torsion}
Let $f : X \to S$ be a smooth proper morphism of schemes.
Let $\Lambda$ be a Noetherian ring. Let $\mathcal{F}$ be a
finite type, locally constant sheaf of $\Lambda$-modules
on $X_\etale$ such that for every geometric point $\overline{x}$ of $X$
the stalk $\mathcal{F}_{\overline{x}}$ is annihilated by an integer
$n > 0$ prime to the residue characteristic of $\overline{x}$.
Then $R^if_*\mathcal{F}$ is a finite type, locally constant sheaf of
$\Lambda$-modules on $S_\etale$ for all $i \in \mathbf{Z}$.
\end{proposition}

\begin{proof}
The question is local on $S$ and hence we may assume $S$ is affine.
For a point $x$ of $X$ denote $n_x \geq 1$
the smallest integer annihilating $\mathcal{F}_{\overline{x}}$
for some (equivalently any) geometric point $\overline{x}$ of $X$
lying over $x$. Since $X$ is quasi-compact (being proper over affine)
there exists a finite \'etale covering $\{U_j \to X\}_{j = 1, \ldots, m}$
such that $\mathcal{F}|_{U_j}$ is constant. Since $U_j \to X$
is open, we conclude that the function $x \mapsto n_x$ is locally
constant and takes finitely many values. Accordingly we
obtain a finite decomposition $X = X_1 \amalg \ldots \amalg X_N$
into open and closed subschemes such that $n_x = n$ if and only if $x \in X_n$.
Then it suffices to prove the lemma for the induced
morphisms $X_n \to S$ and the restriction of $\mathcal{F}$ to $X_n$.
Thus we may and do assume there exists an integer $n > 0$ such
that $\mathcal{F}$ is annihilated by $n$ and such that $n$
is prime to the residue characteristics of all residue fields of $X$.

\medskip\noindent
Since $f$ is smooth and proper the image $f(X) \subset S$ is open and closed.
Hence we may replace $S$ by $f(X)$ and assume $f(X) = S$. In particular,
we see that we may assume $n$ is invertible in the ring defining the
affine scheme $S$.

\medskip\noindent
In this paragraph we reduce to the case where $S$ is Noetherian.
Write $S = \Spec(A)$ for some $\mathbf{Z}[1/n]$-algebra $A$.
Write $A = \bigcup A_i$ as the union of its finite type
$\mathbf{Z}[1/n]$-subalgebras. We can find an $i$ and a morphism
$f_i : X_i \to S_i = \Spec(A_i)$ of finite
type whose base change to $S$ is $f : X \to S$, see
Limits, Lemma \ref{limits-lemma-descend-finite-presentation}.
After increasing $i$ we may assume $f_i : X_i \to S_i$ is
smooth and proper, see Limits, Lemmas
\ref{limits-lemma-eventually-proper}
\ref{limits-lemma-descend-smooth}, and
\ref{limits-lemma-descend-dimension-d}.
By \'Etale Cohomology, Lemma
\ref{etale-cohomology-lemma-category-loc-cst-is-colimit}
we see that there exists an $i$ and a finite type,
locally constant sheaf of $\Lambda$-modules
$\mathcal{F}_i$ whose pullback to $X$ is isomorphic to $\mathcal{F}$.
As $\mathcal{F}$ is annihilated by $n$, we may replace
$\mathcal{F}_i$ by $\Ker(n : \mathcal{F}_i \to \mathcal{F}_i)$
and assume the same thing is true for $\mathcal{F}_i$.
This reduces us to the case discussed in the next paragraph.

\medskip\noindent
Assume we have an integer $n \geq 1$, the base scheme $S$ is Noetherian
and lives over $\mathbf{Z}[1/n]$, and $\mathcal{F}$ is $n$-torsion.
By Theorem \ref{theorem-constructible-shriek}
the sheaves $R^if_*\mathcal{F}$ are constructible sheaves of $\Lambda$-modules.
By \'Etale Cohomology, Lemma
\ref{etale-cohomology-lemma-sp-isom-proper-torsion-loc-ac}
the specialization maps of $R^if_*\mathcal{F}$ are always isomorphisms.
We conclude by \'Etale Cohomology, Lemma
\ref{etale-cohomology-lemma-characterize-locally-constant-module}.
\end{proof}

\end{document}
